\chapter{Cơ sở lý thuyết và yêu cầu hệ thống}
\label{ch:background-requirements}

% Chương hợp nhất (bố cục 6 chương theo quy định của khoa): gộp Ch.2 (Cơ sở lý thuyết)
% và Ch.3 (Yêu cầu hệ thống) cũ thành một chương. Cơ sở lý thuyết đặt trước; phần
% yêu cầu/đặc tả được gộp vào làm các mục cuối.
%
% MIPI-compliance gate (2026-06-24): mọi phát biểu về giao thức bên dưới đã được đối chiếu
% với đặc tả chính thức MIPI I3C Basic v1.1.1 (docs/mipi_i3c_spec.*). Trần SDR 12.5 MHz,
% các tối thiểu thời gian Bảng 87 (24/24/3/12 ns), broadcast 7'h7E, ba điều kiện bus, và các
% mã lệnh ENEC/DISEC/ENTDAA đều khớp.

Chương này thiết lập nền tảng về giao thức và phương pháp luận mà phần còn lại của khóa
luận dựa trên đó, rồi gộp những nền tảng ấy thành các yêu cầu cụ thể, kiểm thử được mà
thiết kế ở \cref{ch:architecture-rtl} thỏa mãn và môi trường kiểm chứng ở
\cref{ch:verif-method} kiểm tra. Phần~A trình bày lý thuyết giao thức --- bus \iic{} mà
thiết kế phải giữ tương thích, chế độ MIPI I3C Basic Single Data Rate (SDR) mà nó hiện
thực, các điều kiện bus và các pha tín hiệu, cơ chế gán địa chỉ động, các Common Command
Code được hỗ trợ, và phương pháp Universal Verification Methodology (UVM) dùng để kiểm
chứng. Phần~B sau đó phát biểu các yêu cầu chức năng, những phần cố ý nằm ngoài phạm vi,
các mục tiêu hiệu năng, giao diện thanh ghi mà phần mềm nhìn thấy, và tiêu chí đóng (closure)
của việc kiểm chứng.

% ===========================================================================
% PHẦN A --- Cơ sở lý thuyết
% ===========================================================================
\section*{Phần A\quad Cơ sở lý thuyết}
\addcontentsline{toc}{section}{Phần A --- Cơ sở lý thuyết}

\section{Tổng quan lại về \iic{}}
\label{sec:bg-i2c}

Bus Inter-Integrated Circuit (\iic{}) là một liên kết nối tiếp hai dây, multi-drop, gồm một
đường xung clock nối tiếp (SCL) và một đường dữ liệu nối tiếp (SDA). Cả hai đường đều ở chế
độ \emph{open-drain} (OD): mỗi thiết bị chỉ có thể kéo đường xuống mức Low, còn một điện trở
pull-up thụ động kéo nó trở lại mức High, nên bus hoạt động như một wired-AND của tất cả các
bộ điều khiển đang lái. Một phiên truyền được đóng khung giữa điều kiện START (SDA bị kéo Low
khi SCL đang High) và điều kiện STOP (SDA được nhả lên High khi SCL đang High). Sau START, bộ
điều khiển dịch ra một byte địa chỉ 9 bit --- 7 bit địa chỉ target theo sau bởi một bit
read/not-write (RnW), bit có trọng số cao nhất đi trước --- và target được địa chỉ trả lời
bằng một acknowledge (ACK) qua việc kéo SDA Low trong xung clock thứ chín; một byte không được
đáp ứng sẽ đọc về thành not-acknowledge (NACK). Dữ liệu sau đó chảy theo từng byte 8 bit, mỗi
byte kèm một ACK/NACK, cho tới khi một STOP đóng khung lại. Một target không theo kịp có thể
giữ SCL ở mức Low (clock stretching) để chặn bộ điều khiển lại.

Ba đặc tính của mô hình này giới hạn \iic{} khi số thiết bị và tốc độ dữ liệu tăng lên, và
I3C được thiết kế để loại bỏ chúng: các pull-up open-drain giới hạn tốc độ clock thực dụng và
lãng phí công suất tĩnh; địa chỉ được cố định khi sản xuất, nên hai linh kiện trùng địa chỉ
không thể cùng tồn tại trên một bus; và bus không có ngắt gốc, buộc bộ điều khiển phải polling.
Những giới hạn này là động lực cho tín hiệu push-pull, gán địa chỉ động, và in-band interrupt
được giới thiệu tiếp sau đây.

\section{MIPI I3C Basic v1.1.1 (SDR, định dạng khung)}
\label{sec:bg-i3c-sdr}

I3C (Improved Inter-Integrated Circuit) giữ lại tô-pô hai dây SCL/SDA và khung START/STOP của
\iic{}, nhưng bổ sung một pha dữ liệu \emph{push-pull} (PP) LV-CMOS \SI{1.8}{\volt} chủ động
lái bus theo cả hai chiều, nâng trần clock Single Data Rate (SDR) lên
\SI{12.5}{\mega\hertz}~\cite{mipi_i3c_basic}. Các pha địa chỉ và phân xử (arbitration) vẫn ở
open-drain để tương thích ngược (\cref{sec:bg-odpp}); chỉ pha dữ liệu chuyển sang push-pull.

I3C dành riêng địa chỉ broadcast \texttt{7'h7E}, mà bộ điều khiển dùng vừa để đánh dấu khởi
đầu hoạt động I3C gốc trên một bus hỗn hợp, vừa để mang các Common Command Code
(\cref{sec:bg-ccc}). Hai dạng khung SDR mang dữ liệu riêng (private) giữa bộ điều khiển và một
target. Trong \textbf{SDR private write}, bộ điều khiển phát một START (có thể đứng trước bởi
broadcast header \texttt{7E}), địa chỉ target với RnW${=}0$, rồi truyền các byte dữ liệu ở pha
push-pull; mỗi byte theo sau bởi một \emph{Transition Bit} (T-Bit) có giá trị là \emph{parity
lẻ} của byte đó, giúp target phát hiện một byte bị hỏng. Trong \textbf{SDR private read}, bộ
điều khiển địa chỉ target với RnW${=}1$ và target lái các byte dữ liệu; lúc này T-Bit mang tín
hiệu \emph{end-of-data} thay vì parity --- target giữ nó ở High khi còn byte tiếp theo và lái
nó về Low ở byte cuối cùng, sau đó bộ điều khiển phát một STOP hoặc Repeated START. Vai trò kép
này của T-Bit (parity khi ghi, end-of-data khi đọc), được tóm tắt trong \cref{fig:sdr-frames},
là trung tâm của các bộ serialiser mức bit được mô tả ở \cref{ch:architecture-rtl}.

\begin{figure}[!htb]
\centering
\begin{tikzpicture}[
  every node/.style={font=\footnotesize},
  field/.style={draw, minimum height=9mm, text height=1.6ex, text depth=.5ex,
                inner xsep=4pt, align=center},
  od/.style={field, fill=blue!10},
  pp/.style={field, fill=orange!18},
  cond/.style={field, fill=black!12, font=\footnotesize\bfseries},
  rowlbl/.style={font=\footnotesize\bfseries, anchor=east},
  drv/.style={font=\scriptsize\itshape, text=black!70},
]

% ---------------- SDR private write (0x7E header enabled) ----------------
\begin{scope}[start chain=w going right, node distance=-\pgflinewidth]
  \node[cond, on chain=w] (w0) at (0,0) {S};
  \node[od,   on chain=w]        {7E\,$+$\,W};
  \node[od,   on chain=w]        {ACK};
  \node[cond, on chain=w]        {Sr};
  \node[pp,   on chain=w]        {addr\,$+$\,W};
  \node[od,   on chain=w]        {ACK};
  \node[pp,   on chain=w]        {Data\,0};
  \node[pp,   on chain=w]        {T};
  \node[pp,   on chain=w]        {$\cdots$};
  \node[pp,   on chain=w]        {Data\,$N\!-\!1$};
  \node[pp,   on chain=w]        {T};
  \node[cond, on chain=w] (wN)   {P};
\end{scope}
\node[rowlbl, left=5pt of w0] {Ghi};
\path (w0.south west) -- (wN.south east)
      node[drv, midway, below=2.5pt] {bộ điều khiển lái các byte dữ liệu \,$\bullet$\, T $=$ parity lẻ của byte};

% ---------------- SDR private read (0x7E header enabled) ----------------
\begin{scope}[start chain=r going right, node distance=-\pgflinewidth]
  \node[cond, on chain=r] (r0) at (0,-2.5) {S};
  \node[od,   on chain=r]        {7E\,$+$\,W};
  \node[od,   on chain=r]        {ACK};
  \node[cond, on chain=r]        {Sr};
  \node[pp,   on chain=r]        {addr\,$+$\,R};
  \node[od,   on chain=r]        {ACK};
  \node[pp,   on chain=r]        {Data\,0};
  \node[pp,   on chain=r]        {T$=$1};
  \node[pp,   on chain=r]        {$\cdots$};
  \node[pp,   on chain=r]        {Data\,$N\!-\!1$};
  \node[pp,   on chain=r]        {T$=$0};
  \node[cond, on chain=r] (rN)   {P};
\end{scope}
\node[rowlbl, left=5pt of r0] {Đọc};
\path (r0.south west) -- (rN.south east)
      node[drv, midway, below=2.5pt] {target lái các byte dữ liệu \,$\bullet$\, T $=$ end-of-data (1 $=$ còn byte, 0 $=$ byte cuối)};

% ---------------- legend ----------------
\begin{scope}[every node/.style={font=\scriptsize}, node distance=3pt]
  \node[cond, minimum width=6mm, minimum height=4mm] (lg0) at (0,-4.35) {};
  \node[right=of lg0] (lg0t) {S\,/\,Sr\,/\,P (điều kiện bus)};
  \node[od, minimum width=6mm, minimum height=4mm, right=12pt of lg0t] (lg1) {};
  \node[right=of lg1] (lg1t) {Open-drain};
  \node[pp, minimum width=6mm, minimum height=4mm, right=12pt of lg1t] (lg2) {};
  \node[right=of lg2] {Push-pull};
\end{scope}
\end{tikzpicture}
\caption{Định dạng khung SDR private write và private read (với broadcast header
\texttt{0x7E} được bật).}
\label{fig:sdr-frames}
\end{figure}

\section{Các điều kiện bus (START / Repeated START / STOP)}
\label{sec:bg-bus-conditions}

Ba điều kiện cạnh (edge) phân định và nối các phiên truyền. Một \textbf{START} (S) là cạnh
xuống của SDA khi SCL đang High; một \textbf{STOP} (P) là cạnh lên của SDA khi SCL đang High;
và một \textbf{Repeated START} (Sr) là một START được phát mà không có STOP xen giữa, dùng để
nối các pha --- ví dụ để chuyển từ broadcast header open-drain sang pha dữ liệu push-pull, hoặc
để đảo chiều bus giữa một lần ghi và một lần đọc.

Đặc tả MIPI còn định nghĩa ba điều kiện rảnh (idle) \emph{khác nhau} mà khóa luận này cẩn thận
không nhập nhằng~\cite{mipi_i3c_basic}. \textbf{Bus Free Condition} (\S5.1.3.2.1) là khoảng sau
một STOP và trước START kế tiếp, kéo dài ít nhất $t_{\mathrm{CAS}}$; \textbf{Bus Available
Condition} (\S5.1.3.2.2) là một Bus Free Condition kéo dài ít nhất $t_{\mathrm{AVAL}}$, trong đó
một target có thể yêu cầu lái bus; và \textbf{Bus Idle Condition} (\S5.1.3.2.3) là một khoảng
dài hơn (ít nhất $t_{\mathrm{IDLE}}$) cần có trước một yêu cầu Hot-Join. Bộ điều khiển này tạo
ra START/Sr/STOP và phát hiện Bus Free Condition trước khi khởi tạo một phiên truyền; nó không
dùng Bus Available hay Bus Idle Condition, vì các sự kiện cần đến chúng (START khởi tạo từ
target, Hot-Join) nằm ngoài phạm vi. Bus monitor và SCL generator hiện thực các điều kiện này
được trình bày chi tiết ở \cref{ch:architecture-rtl}.

\figph{Các điều kiện bus}{START / Repeated START (Sr) / STOP --- cạnh SDA so với SCL High}
  {MIPI I3C Basic v1.1.1}{WaveDrom (bus/bus_conditions.json)}

\section{Tín hiệu open-drain so với push-pull}
\label{sec:bg-odpp}

I3C dùng hai chế độ tín hiệu điện trên cùng các dây và chuyển qua lại giữa chúng ngay trong một
phiên truyền. \textbf{Các bit acknowledge và các pha địa chỉ có phân xử giữ ở open-drain}: bus
được nhả ra để bị kéo lên High và bất kỳ thiết bị nào cũng có thể kéo nó Low, điều này giữ lại
phân xử kiểu wired-AND của \iic{}, cho phép các target \iic{} hợp pháp tham gia, và làm cho phân
xử ENTDAA theo từng bit ở \cref{sec:bg-entdaa} khả thi. Một khi bộ điều khiển đã thắng bus,
\textbf{pha dữ liệu --- và mọi địa chỉ mà nó lái sau một Repeated START --- chuyển sang
push-pull}, nơi bộ lái đang hoạt động cấp và rút dòng trực tiếp; điều này loại bỏ hằng số thời
gian RC của pull-up và chính là thứ cho phép tốc độ SDR \SI{12.5}{\mega\hertz}.

Ranh giới được vượt qua tại một Repeated START: broadcast header \texttt{0x7E} và ACK của nó
được gửi open-drain, và tại Sr bộ điều khiển xoay bus sang push-pull, nên địa chỉ target theo
sau và tất cả các byte dữ liệu được lái push-pull trong khi các bit ACK/NACK vẫn ở open-drain
(\cref{fig:sdr-frames}). Khi broadcast header bị tắt, bộ điều khiển thay vào đó gửi địa chỉ
target ở open-drain ngay sau START --- không có Repeated START --- và chỉ chuyển sang push-pull
cho pha dữ liệu. Trong RTL, lựa chọn chế độ này được phơi bày như một tín hiệu duy nhất,
\texttt{sel_od_pp} (\texttt{1}${=}$push-pull, \texttt{0}${=}$open-drain), được lái vào lớp PHY
(\cref{ch:architecture-rtl}).

\figph{Các pha open-drain so với push-pull}
  {OD 7E-header/ACK so với PP địa chỉ+dữ liệu; ranh giới OD$\rightarrow$PP tại Sr; lựa chọn \texttt{sel_od_pp}}
  {MIPI I3C Basic v1.1.1}{WaveDrom (bus/od_vs_pp.json)}

\section{Gán địa chỉ động qua ENTDAA}
\label{sec:bg-entdaa}

I3C thay cơ chế địa chỉ cố định của \iic{} bằng Dynamic Address Assignment (DAA) lúc chạy, qua
đó loại bỏ va chạm địa chỉ và hỗ trợ các thiết bị dùng chung mã linh kiện. Thiết kế này hiện
thực DAA thông qua thủ tục ENTDAA (Enter Dynamic Address Assignment) từ góc nhìn của bộ điều
khiển. Bộ điều khiển phát một START theo sau bởi địa chỉ broadcast \texttt{7'h7E}${+}$W và
Common Command Code ENTDAA (\texttt{0x07}); sau đó nó vào một vòng lặp, phát một Repeated START
và \texttt{7'h7E}${+}$R cho mỗi vòng. Mỗi target chưa được gán địa chỉ trả lời bằng cách dòng
ra danh tính 64 bit duy nhất của nó --- một Provisioned ID (PID) 48 bit, một Bus Characteristics
Register (BCR) 8 bit, và một Device Characteristics Register (DCR) 8 bit --- bit trọng số cao
nhất đi trước, không có ACK giữa các trường.

Vì dòng danh tính ở open-drain và wired-AND, nó đồng thời đóng vai trò cơ chế phân xử: khi nhiều
target đáp ứng cùng lúc, một target lái logic~1 nhưng đọc về~0 trên SDA tức là đã thua một danh
tính có giá trị nhỏ hơn và rút lui, nên đúng một target thắng mỗi vòng. Bộ điều khiển sau đó gửi
cho người thắng một địa chỉ động 7 bit kèm một bit parity, target ACK, và vòng lặp lặp lại cho
tới khi một vòng không có đáp ứng, lúc đó một STOP đóng thủ tục lại. Các địa chỉ đã gán được ghi
vào một Device Address Table (DAT) 32 mục. \Cref{ch:architecture-rtl} hiện thực điều này thành
một phân hệ hai FSM --- một bộ quản lý vòng lặp và một FSM phân xử theo từng vòng.

\figph{Phân xử ENTDAA}
  {7E+W, ACK, ENTDAA CCC, rồi theo từng target: Sr, 7E+R, PID+BCR+DCR 64-bit, dyn-addr+parity,
   ACK; lặp tới khi không còn đáp ứng, rồi P}
  {MIPI I3C Basic v1.1.1}{WaveDrom (bus/entdaa_arbitration.json)}

\section{Tổng quan Common Command Code (CCC)}
\label{sec:bg-ccc}

Common Command Code (CCC) là các lệnh chuẩn hóa mà bộ điều khiển phát sau broadcast header
\texttt{7'h7E} để cấu hình và quản lý các target. Chúng có hai dạng khung. Một \textbf{Broadcast
CCC} gửi tới tất cả các target cùng lúc (\texttt{[S]} \texttt{[7E+W]} \texttt{[CCC]} \dots\
\texttt{[P]}); một \textbf{Direct CCC} nhắm tới một thiết bị bằng cách theo sau mã lệnh là một
Repeated START và địa chỉ của thiết bị đó (\texttt{[S]} \texttt{[7E+W]} \texttt{[CCC]}
\texttt{[Sr]} \texttt{[addr+RnW]} \dots\ \texttt{[P]}).

Trong toàn bộ danh mục CCC, thiết kế này chỉ hiện thực tập con thiết yếu cần để khởi động và
quản lý một bus SDR, được tóm tắt trong \cref{tab:ccc-supported}. ENTDAA (\texttt{0x07},
broadcast) điều khiển việc gán địa chỉ động; ENEC (\texttt{0x00} broadcast / \texttt{0x80}
direct) và DISEC (\texttt{0x01} broadcast / \texttt{0x81} direct) bật và tắt các sự kiện do
target tạo ra. Mọi CCC khác đều nằm ngoài phạm vi (\cref{sec:req-oos}).

\begin{table}[!htb]\centering
  \caption{Các Common Command Code được hiện thực trong thiết kế này.}
  \label{tab:ccc-supported}
  \small
  \begin{tabular}{@{}llcl@{}}
    \toprule
    \textbf{CCC} & \textbf{Dạng khung} & \textbf{Mã} & \textbf{Mục đích} \\
    \midrule
    ENTDAA & Broadcast        & \texttt{0x07} & Vào chế độ gán địa chỉ động \\
    ENEC   & Broadcast        & \texttt{0x00} & Bật sự kiện target (tất cả target) \\
    DISEC  & Broadcast        & \texttt{0x01} & Tắt sự kiện target (tất cả target) \\
    ENEC   & Direct           & \texttt{0x80} & Bật sự kiện target (một target) \\
    DISEC  & Direct           & \texttt{0x81} & Tắt sự kiện target (một target) \\
    \bottomrule
  \end{tabular}
\end{table}

\figph{Định dạng khung CCC}{Khung Broadcast CCC so với Direct CCC}
  {MIPI I3C Basic v1.1.1}{WaveDrom (bus/ccc_frame.json)}

\section{So sánh I3C với \iic{}}
\label{sec:bg-comparison}

\Cref{tab:i3c-vs-i2c} tổng hợp các khác biệt giới thiệu ở trên thành một ma trận tính năng và
đánh dấu những hàng thuộc phạm vi của thiết kế này. Các lợi ích nổi bật là tốc độ dữ liệu tăng
$\sim$12.5$\times$ nhờ tín hiệu push-pull, gán địa chỉ động không va chạm, và công suất tĩnh
thấp hơn do không có dòng pull-up ở mức bo mạch --- tất cả trong khi vẫn tương thích ngược với
lưu lượng \iic{} Fast-Mode hợp pháp trên cùng hai dây. \Cref{tab:timing-params} sau đó liệt kê
các tham số thời gian SDR push-pull và \iic{} Fast-Mode mà phần định thời lập trình được của bộ
điều khiển phải thỏa mãn; các giá trị SDR lấy từ Bảng~87 của MIPI I3C Basic và các giá trị
Fast-Mode lấy từ bảng định thời \iic{} hợp pháp~\cite{mipi_i3c_basic}.

\begin{table}[!htb]\centering
  \caption{Ma trận tính năng I3C so với \iic{}. Cột cuối đánh dấu mức độ bao phủ trong thiết kế này.}
  \label{tab:i3c-vs-i2c}
  \small
  \begin{tabular}{@{}>{\raggedright\arraybackslash}p{0.21\textwidth}
                     >{\raggedright\arraybackslash}p{0.27\textwidth}
                     >{\raggedright\arraybackslash}p{0.30\textwidth}
                     >{\raggedright\arraybackslash}p{0.12\textwidth}@{}}
    \toprule
    \textbf{Tính năng} & \textbf{\iic{} (Fm / Fm+)} & \textbf{I3C SDR} & \textbf{Trong phạm vi} \\
    \midrule
    Tần số SCL tối đa   & \SI{400}{\kilo\hertz} (Fm) / \SI{1}{\mega\hertz} (Fm+) & \SI{12.5}{\mega\hertz} (push-pull) & Có (Fm) \\
    Tín hiệu            & Chỉ open-drain             & Open-drain (addr/arb) + push-pull (data) & Có \\
    Địa chỉ             & 7-bit tĩnh                 & 7-bit động               & Có \\
    Cách gán địa chỉ    & Cố định / cấu hình phần mềm & Động (ENTDAA)            & Có \\
    Pull-up bus         & Điện trở ngoài trên bo     & Bộ điều khiển cung cấp; không cần điện trở ngoài & Có \\
    Clock stretching    & Target có thể kéo dài SCL  & Target không thể; bộ điều khiển có thể chặn clock & Không áp dụng \\
    Phát hiện lỗi       & Chỉ ACK/NACK               & ACK/NACK + parity T-Bit ($+$ CRC ở HDR) & Parity có; CRC không \\
    In-Band Interrupt   & Không (polling)            & Hỗ trợ IBI               & Không \\
    Hot-Join            & Không                      & Hỗ trợ                   & Không \\
    CCC                 & Không                      & Tập lệnh Broadcast / Direct & Chỉ tập con \\
    Tương thích ngược   & ---                        & \iic{} hợp pháp trên cùng bus & Có \\
    Số dây              & 2 (SCL + SDA)              & 2 (SCL + SDA)            & Có \\
    \bottomrule
  \end{tabular}
\end{table}

\begin{table}[!htb]\centering
  \caption{Các tham số thời gian SDR push-pull và \iic{} Fast-Mode (giá trị tối thiểu trừ khi ghi chú khác).}
  \label{tab:timing-params}
  \small
  \begin{tabular}{@{}>{\raggedright\arraybackslash}p{0.34\textwidth} l r r l@{}}
    \toprule
    \textbf{Tham số} & \textbf{Ký hiệu} & \textbf{I3C SDR} & \textbf{\iic{}-Fm} & \textbf{Đơn vị} \\
    \midrule
    Tần số SCL (tối đa)            & $f_{\mathrm{SCL}}$  & 12.5 (điển hình) & 400  & \si{\kilo\hertz}/\si{\mega\hertz} \\
    Chu kỳ SCL Low                 & $t_{\mathrm{LOW}}$  & 24   & 1300 & \si{\nano\second} \\
    Chu kỳ SCL High                & $t_{\mathrm{HIGH}}$ & 24   & 600  & \si{\nano\second} \\
    Thời gian setup dữ liệu        & $t_{\mathrm{SU,DAT}}$ & 3  & 100  & \si{\nano\second} \\
    Thời gian hold dữ liệu         & $t_{\mathrm{HD,DAT}}$ & 0  & 0    & \si{\nano\second} \\
    Setup START                    & $t_{\mathrm{SU,STA}}$ & --- & 600 & \si{\nano\second} \\
    Hold START                     & $t_{\mathrm{HD,STA}}$ & --- & 600 & \si{\nano\second} \\
    Setup STOP                     & $t_{\mathrm{SU,STO}}$ & 12 & ---  & \si{\nano\second} \\
    Clock-to-data ra (target, tối đa)& $t_{\mathrm{SCO}}$  & 12   & ---  & \si{\nano\second} \\
    Bus free                       & $t_{\mathrm{CAS}}$/$t_{\mathrm{BUF}}$ & 38 & 1300 & \si{\nano\second} \\
    \bottomrule
  \end{tabular}
\end{table}

\section{Tổng quan phương pháp UVM 1.2}
\label{sec:bg-uvm}

Thiết kế được kiểm chứng trong một testbench dựng bằng Universal Verification Methodology (UVM)
1.2, một thư viện lớp SystemVerilog chuẩn hóa để xây dựng các môi trường kiểm chứng phân lớp,
tái sử dụng được~\cite{uvm12}. UVM tổ chức một testbench thành một cây phân cấp các thành phần.
Một \texttt{uvm\_test} chọn kịch bản để chạy (chọn lúc chạy qua \texttt{+UVM\_TESTNAME}, nên các
kịch bản không cần biên dịch lại) và tạo một \texttt{uvm\_env}, chứa hạ tầng kiểm chứng tái sử
dụng được: một hoặc nhiều \emph{agent}, một scoreboard, và các bộ thu coverage. Mỗi
\texttt{uvm\_agent} gói ba thành phần con --- một \emph{driver} biến các giao dịch trừu tượng
thành hoạt động mức chân (pin) trên một virtual interface, một \emph{monitor} quan sát các chân
và tái dựng các giao dịch, và một \emph{sequencer} điều phối dòng các phần tử kích thích
(stimulus) từ các sequence tới driver.

Kích thích được mô tả tách rời với cấu trúc: một \texttt{uvm\_sequence} sinh ra các giao dịch
\texttt{uvm\_sequence\_item} và phát chúng qua một sequencer, nên cùng một môi trường có thể chạy
nhiều kịch bản khác nhau. Các thành phần giao tiếp qua các analysis port theo \emph{Transaction-Level
Modelling} (TLM) --- một monitor phát quảng bá mỗi giao dịch bắt được trên một
\texttt{uvm\_analysis\_port}, port này tỏa ra tới scoreboard và các bộ thu coverage mà bên sản
xuất không cần biết bên tiêu thụ là ai. Hai dịch vụ nữa gắn kết môi trường lại: \texttt{uvm\_config\_db}
truyền cấu hình (các handle virtual-interface, các knob) xuống cây phân cấp theo tên phân cấp, và
\emph{factory} (\texttt{type\_id::create}) tạo các thành phần và giao dịch một cách gián tiếp để
có thể ghi đè kiểu mà không phải sửa testbench. Các cơ chế này là nền tảng cho môi trường hai
agent được mô tả ở \cref{ch:verif-method}.

\figph{Phân lớp testbench UVM}
  {test $\rightarrow$ env $\rightarrow$ agent (driver/monitor/sequencer); các analysis port TLM
   tới scoreboard/coverage; config\_db và factory}{Accellera UVM 1.2}{TikZ}

% ===========================================================================
% PHẦN B --- Yêu cầu và đặc tả hệ thống
% ===========================================================================
\section*{Phần B\quad Yêu cầu và đặc tả hệ thống}
\addcontentsline{toc}{section}{Phần B --- Yêu cầu và đặc tả hệ thống}

\section{Yêu cầu chức năng}
\label{sec:req-functional}

Các hành vi mà bộ điều khiển phải hiện thực được liệt kê thành các hạng mục kiểm thử được trong
\cref{tab:func-req}; mỗi hạng mục được thực thi bởi ít nhất một virtual sequence có chủ đích ở
\cref{ch:verif-method} và được scoreboard đối chiếu chéo. Danh sách trải khắp toàn bộ đường dữ
liệu SDR (private write, read, và immediate transfer), khả năng phối hợp với \iic{} Fast-Mode hợp
pháp, cơ chế điều kiện bus và tín hiệu, gán địa chỉ động, các CCC được hỗ trợ, và đường báo cáo
mà phần mềm nhìn thấy.

\begin{table}[!htb]\centering
  \caption{Yêu cầu chức năng.}
  \label{tab:func-req}
  \small
  \begin{tabular}{@{}l >{\raggedright\arraybackslash}p{0.78\textwidth}@{}}
    \toprule
    \textbf{ID} & \textbf{Yêu cầu} \\
    \midrule
    FR-1  & SDR private write: pha dữ liệu push-pull với parity lẻ T-Bit theo từng byte \\
    FR-2  & SDR private read: dữ liệu do target lái với tín hiệu end-of-data của T-Bit \\
    FR-3  & Immediate Data Transfer: tối đa bốn byte inline mang trong command descriptor \\
    FR-4  & \iic{} Fast-Mode legacy write (\SI{400}{\kilo\hertz}, định thời cố định) \\
    FR-5  & \iic{} Fast-Mode legacy read \\
    FR-6  & Tạo các điều kiện START, Repeated START (Sr), và STOP \\
    FR-7  & Phát hiện Bus Free Condition trước khi khởi tạo một phiên truyền \\
    FR-8  & Chuyển pha open-drain$\leftrightarrow$push-pull qua \texttt{sel_od_pp} \\
    FR-9  & Vòng lặp DAA đa target ENTDAA, hậu thuẫn bởi Device Address Table (DAT) 32 mục \\
    FR-10 & ENEC / DISEC ở cả hai dạng khung broadcast và direct \\
    FR-11 & Hợp đồng đọc/ghi thanh ghi 32-bit (staging lệnh, buffer dữ liệu) \\
    FR-12 & Báo cáo Response Queue (mã trạng thái lỗi và độ dài dữ liệu cho mỗi lệnh) \\
    FR-13 & RX Data Buffer báo cáo dữ liệu đọc nhận được tới phần mềm host \\
    \bottomrule
  \end{tabular}
\end{table}

\section{Các tính năng ngoài phạm vi}
\label{sec:req-oos}

Để giữ thiết kế khả thi trong một đồ án 10 tín chỉ, một vài lát cắt độc lập của đặc tả I3C được
cố ý loại trừ. Mỗi loại trừ trong \cref{tab:oos} bỏ đi một tính năng không cần để chứng minh một
đường dữ liệu SDR đầy đủ, chính xác cùng khả năng phối hợp với thiết bị hợp pháp; một số được bàn
lại như hướng phát triển ở \cref{ch:conclusion}. Đây là các quyết định thiết kế, nhất quán với
phạm vi nêu ở \cref{ch:intro-full}.

\begin{table}[!htb]\centering
  \caption{Các tính năng ngoài phạm vi và lý do.}
  \label{tab:oos}
  \small
  \begin{tabular}{@{}>{\raggedright\arraybackslash}p{0.34\textwidth}
                     >{\raggedright\arraybackslash}p{0.58\textwidth}@{}}
    \toprule
    \textbf{Tính năng} & \textbf{Lý do loại trừ} \\
    \midrule
    In-Band Interrupt (IBI)            & Tính năng báo sự kiện; vận hành không-polling không cần để chứng minh đường dữ liệu SDR \\
    Hot-Join                           & Cắm thiết bị lúc chạy; cần Bus Idle Condition, không dùng ở đây \\
    Các chế độ HDR (DDR / TSP / TSL)   & Các chế độ tốc độ cao hơn ngoài chế độ SDR bắt buộc \\
    \iic{} Fast-Mode Plus (\SI{1}{\mega\hertz}) & Fast-Mode (\SI{400}{\kilo\hertz}) đã đủ cho khả năng phối hợp hợp pháp \\
    Secondary / multi-controller       & Chỉ vai trò Active-Controller đơn; không có chuyển giao bộ điều khiển \\
    Chế độ Target (slave)              & Thiết kế chỉ làm controller \\
    Bus recovery                       & Leo thang khôi phục lỗi vượt quá phần báo lỗi trong phạm vi \\
    Tuân thủ HCI đầy đủ / DCT          & Thay bằng giao diện thanh ghi 32-bit đơn giản hóa \\
    Mọi CCC khác                       & Chỉ ENTDAA, ENEC, và DISEC được hiện thực \\
    \bottomrule
  \end{tabular}
\end{table}

\section{Mục tiêu hiệu năng}
\label{sec:req-performance}

Bộ điều khiển nhắm tới SCL push-pull SDR tối đa \SI{12.5}{\mega\hertz} và SCL \iic{} Fast-Mode
hợp pháp \SI{400}{\kilo\hertz}, cả hai đều giữ ở mức hoặc dưới các trần đặc tả của
\cref{tab:timing-params}. Vì SCL được tạo bằng cách đếm chu kỳ clock hệ thống, thiết kế chạy
trong một miền clock đơn ở tối thiểu \SI{333}{\mega\hertz} (chu kỳ \SI{3}{\nano\second};
testbench cấp clock cho DUT đúng tốc độ này). Tốc độ này được đặt bởi ngân sách clock-to-data-out
của target: $t_{\mathrm{SCO,max}} = 4 \times T_{\mathrm{sys}} = 4 \times \SI{3}{\nano\second} =
\SI{12}{\nano\second}$, đáp ứng giá trị tối đa \SI{12}{\nano\second} của \cref{tab:timing-params}.

Định thời SCL không cố định cứng mà lập trình được qua các control/status register (CSR); pha
Low và High của SCL được nạp từ các thanh ghi đếm theo từng chế độ, nên phần mềm có thể tinh
chỉnh bus mà không cần đổi RTL. \Cref{tab:perf-targets} liệt kê các giá trị đếm mặc định khi reset
lấy trực tiếp từ register file của RTL. Có các giá trị đếm chu kỳ Low riêng cho pha dữ liệu
push-pull và pha địa chỉ open-drain chậm hơn ($t_{\mathrm{LOW,OD}}$), và một tập đếm riêng chi
phối \iic{} Fast-Mode. Các mặc định được chọn để nằm trong các trần đặc tả với biên dự phòng;
tần số đạt được suy ra từ các giá trị đếm đã nạp và chu kỳ clock hệ thống.

\begin{table}[!htb]\centering
  \caption{Mục tiêu hiệu năng và giá trị đếm CSR định thời mặc định khi reset (theo chu kỳ
  clock hệ thống ở \SI{3}{\nano\second}).}
  \label{tab:perf-targets}
  \small
  \begin{tabular}{@{}>{\raggedright\arraybackslash}p{0.40\textwidth} l r@{}}
    \toprule
    \textbf{Mục tiêu / tham số} & \textbf{CSR} & \textbf{Mặc định} \\
    \midrule
    Clock hệ thống (miền đơn)            & ---              & $\ge$\SI{333}{\mega\hertz} \\
    Trần SCL SDR                        & ---              & $\le$\SI{12.5}{\mega\hertz} \\
    Trần SCL \iic{}-Fm                  & ---              & \SI{400}{\kilo\hertz} \\
    \addlinespace
    SCL Low SDR (push-pull)             & \texttt{T\_LOW}     & 16 \\
    SCL Low SDR (pha địa chỉ open-drain)& \texttt{T\_LOW\_OD} & 67 \\
    SCL High SDR                        & \texttt{T\_HIGH}    & 11 \\
    Rise / fall SDR                     & \texttt{T\_R} / \texttt{T\_F} & 4 / 4 \\
    SCL Low / High \iic{}-Fm            & \texttt{I2C\_T\_LOW} / \texttt{I2C\_T\_HIGH} & 534 / 300 \\
    Rise / fall \iic{}-Fm               & \texttt{I2C\_T\_R} / \texttt{I2C\_T\_F} & 100 / 100 \\
    \addlinespace
    Số mục Device Address Table         & \texttt{DAT\_DEPTH}  & 32 \\
    Độ sâu FIFO CMD/TX/RX/RESP (mỗi cái) & \texttt{*\_FIFO\_DEPTH} & 64 \\
    \bottomrule
  \end{tabular}
\end{table}

\section{Giao diện phần mềm nhìn thấy (bus thanh ghi 32-bit)}
\label{sec:req-swif}

Bộ điều khiển được phơi bày cho phần mềm host qua một giao diện thanh ghi 32-bit cố ý đơn giản
--- không phải AXI4 hay AHB-Lite như trong thiết kế tham chiếu --- được chọn để tập trung vào
hành vi giao thức hơn là logic bộ chuyển đổi bus. Một lần ghi đưa ra một địa chỉ, dữ liệu ghi, và
một write-enable trong một chu kỳ (\texttt{reg_addr_i[11:0]} ${+}$ \texttt{reg_wdata_i[31:0]}
${+}$ \texttt{reg_wen_i}); một lần đọc đưa ra một địa chỉ và một read-enable rồi nhận dữ liệu
kèm một bắt tay ready (\texttt{reg_addr_i[11:0]} ${+}$ \texttt{reg_ren_i} $\rightarrow$
\texttt{reg_rdata_o[31:0]} ${+}$ \texttt{reg_ready_o}).

Bên trên hợp đồng thanh ghi này là các cấu trúc staging giao dịch: một lệnh được staging thành
một command descriptor hai DWORD trước khi được đẩy vào command queue; payload ghi được viết qua
một cửa sổ TX data-buffer và payload đọc được lấy qua một cửa sổ RX data-buffer; và mỗi lệnh hoàn
tất tạo ra một response descriptor trong response queue. Bản đồ thanh ghi đầy đủ và các bố cục
descriptor được dời sang phần phụ lục và trình bày chi tiết ở \cref{ch:architecture-rtl}.

\figph{Giao thức đọc/ghi bus thanh ghi}
  {bắt tay ghi một chu kỳ + định thời đọc-kèm-ready trên bus thanh ghi 32-bit}
  {csr\_registers.sv}{WaveDrom / TikZ timing}

\section{Yêu cầu kiểm chứng}
\label{sec:req-verification}

Việc kiểm chứng được xem là đóng (closed) khi: mọi virtual sequence có chủ đích trong danh mục
chức năng của nó đều pass; các đối chiếu chéo của scoreboard về trạng thái command, response,
receive-data, CSR, và DAT đều sạch; các SystemVerilog Assertion (SVA) gắn vào testbench vẫn im
lặng; và các mục tiêu regression được phân loại chạy xanh. Một phần then chốt của hợp đồng này là
scoreboard kiểm tra bộ điều khiển so với \emph{đúng} tập mã trạng thái lỗi mà thiết kế được yêu
cầu tạo ra --- không nhiều hơn cũng không ít hơn.

\Cref{tab:err-codes} liệt kê cách mã hóa error-status của response. Bảng liệt kê khớp với định
nghĩa HCI/TCRI để tương thích, nhưng chỉ những mã được đánh dấu \emph{được tạo} là do thiết kế
chỉ-SDR này phát ra và do đó nằm trong hợp đồng kiểm chứng. Các mã chỉ-HDR (\texttt{Crc},
\texttt{Parity}, \texttt{Frame}) được khai báo nhưng không bao giờ được tạo, vì các chế độ HDR
nằm ngoài phạm vi. Hai mã cần lưu ý vì vai trò của chúng dễ bị hoán đổi: \texttt{Nack}
(\texttt{0x5}) dành cho đường ENTDAA khi một địa chỉ động bị từ chối hai lần (hết lượt retry),
\emph{không phải} cho một NACK dữ liệu; mã thực sự được trả về cho một NACK pha dữ liệu hoặc một
bus abort là \texttt{I2cDataNackOrI3cBusAborted} (\texttt{0x9}). Phân biệt này chính là điều mà
logic ánh xạ response của command FSM thực thi ở \cref{ch:architecture-rtl}.

\begin{table}[!htb]\centering
  \caption{Các mã error-status của response. Chỉ các mã \emph{được tạo} nằm trong hợp đồng
  kiểm chứng.}
  \label{tab:err-codes}
  \small
  \begin{tabular}{@{}c l >{\raggedright\arraybackslash}p{0.44\textwidth} c@{}}
    \toprule
    \textbf{Mã} & \textbf{Tên} & \textbf{Ý nghĩa} & \textbf{Được tạo} \\
    \midrule
    \texttt{0x0} & \texttt{Success}                    & Không lỗi                                       & Có \\
    \texttt{0x1} & \texttt{Crc}                        & Lỗi CRC (chế độ HDR)                            & Không \\
    \texttt{0x2} & \texttt{Parity}                     & Sai parity T-Bit (chế độ HDR)                   & Không \\
    \texttt{0x3} & \texttt{Frame}                      & Vi phạm khung (chế độ HDR)                      & Không \\
    \texttt{0x4} & \texttt{AddrHeader}                 & NACK address-header                             & Có \\
    \texttt{0x5} & \texttt{Nack}                       & Địa chỉ DAA bị từ chối hai lần (hết retry ENTDAA) & Có \\
    \texttt{0x6} & \texttt{Ovl}                        & Tràn / cạn FIFO                                 & Có \\
    \texttt{0x7} & \texttt{I3cShortReadErr}            & Target kết thúc đọc sớm (T-Bit${=}0$ trước độ dài) & Có \\
    \texttt{0x8} & \texttt{HcAborted}                  & Host-controller yêu cầu abort                   & Có \\
    \texttt{0x9} & \texttt{I2cDataNackOrI3cBusAborted} & NACK pha dữ liệu \iic{} hoặc abort bus I3C       & Có \\
    \texttt{0xA} & \texttt{NotSupported}               & Loại lệnh không hỗ trợ                          & Có \\
    \bottomrule
  \end{tabular}
\end{table}
